\chapter*{缩写与术语表}
\addcontentsline{toc}{chapter}{缩写与术语表}

本表只收录本册反复出现的缩写。缩写不是背诵目标；它们的价值是帮助读者把较长的机制名称
与正文中的状态对象和协议边界对应起来。

\begin{longtable}{|L{0.22\linewidth}|L{0.31\linewidth}|L{0.41\linewidth}|}
\toprule
缩写 & 全称或常见含义 & 本册语境 \\
\midrule
ABI & Application Binary Interface & 系统调用参数、调用约定、用户栈、ELF 装载契约 \\\tablerowrule
ACPI & Advanced Configuration and Power Interface & 固件向内核描述 CPU、NUMA、设备和电源信息 \\\tablerowrule
ALU & Arithmetic Logic Unit & CPU 数据通路中执行整数算术和逻辑运算的部件 \\\tablerowrule
APIC & Advanced Programmable Interrupt Controller & x86 中断投递、本地定时器、IPI 和多核启动 \\\tablerowrule
PCID & Process Context ID & x86-64 TLB 中区分地址空间的标签，减少切换时全局刷新 \\\tablerowrule
BIO & Block I/O & 块层从文件系统接收的页片段读写意图 \\\tablerowrule
BPF & Berkeley Packet Filter & seccomp、eBPF tracing、网络和安全钩子的执行载体 \\\tablerowrule
CFS & Completely Fair Scheduler & Linux 经典公平调度模型，以 \code{vruntime} 解释按权重分配 CPU \\\tablerowrule
CFI & Control Flow Integrity & 控制流保护，限制间接跳转和返回路径被劫持 \\\tablerowrule
CPL & Current Privilege Level & x86-64 当前特权级，决定是否可执行特权操作 \\\tablerowrule
COW & Copy-on-Write & \code{fork}、\code{MAP\_PRIVATE} 和 COW 文件系统的共享后复制语义 \\\tablerowrule
CR/MSR & Control Register / Model-Specific Register & x86-64 控制寄存器和模型专用寄存器，保存页表根、系统调用入口和特权控制状态 \\\tablerowrule
DMA & Direct Memory Access & 设备直接读写内存，要求映射、缓存一致性和生命周期管理 \\\tablerowrule
EEVDF & Earliest Eligible Virtual Deadline First & Linux 自 6.6 起逐步采用的公平调度策略，结合 lag 与虚拟截止期选择任务 \\\tablerowrule
EOI & End Of Interrupt & 中断处理后通知控制器可以投递后续中断 \\\tablerowrule
FUA & Force Unit Access & 块设备写入要求到达稳定介质，而非仅进入设备缓存 \\\tablerowrule
GDT & Global Descriptor Table & x86-64 早期保护和段描述结构，长模式下仍参与入口设置 \\\tablerowrule
GFP & Get Free Page flags & 内核分配上下文，决定是否可睡眠、可 I/O、可回收 \\\tablerowrule
IDT & Interrupt Descriptor Table & x86-64 中断、异常和陷入入口表 \\\tablerowrule
IOMMU & I/O Memory Management Unit & 设备 DMA 地址翻译和隔离 \\\tablerowrule
IPI & Inter-Processor Interrupt & CPU 间重调度、TLB shootdown、多核控制消息 \\\tablerowrule
IRQ & Interrupt Request & 外设中断请求及其内核处理路径 \\\tablerowrule
ISA & Instruction Set Architecture & 软件可见的指令、寄存器、异常、内存和特权契约 \\\tablerowrule
LSM & Linux Security Module & SELinux、AppArmor、BPF LSM 等安全钩子框架 \\\tablerowrule
MMU & Memory Management Unit & 执行虚拟地址翻译和页表权限检查的硬件单元 \\\tablerowrule
MMIO & Memory-Mapped I/O & CPU 通过地址空间读写设备寄存器 \\\tablerowrule
NAPI & New API & Linux 网络收包中断转轮询机制，控制中断风暴 \\\tablerowrule
NUMA & Non-Uniform Memory Access & 多插槽内存局部性、调度和页分配策略 \\\tablerowrule
OOM & Out Of Memory & 内存耗尽、回收失败和 OOM killer 决策 \\\tablerowrule
PCIe & Peripheral Component Interconnect Express & 现代外设互联，承载配置空间、BAR、MSI/MSI-X 和 DMA \\\tablerowrule
PTE/PMD/PUD/PGD & Page Table Entry 等页表层级 & 多级页表、权限位、缺页和 TLB 失效 \\\tablerowrule
RCU & Read-Copy-Update & 读路径低成本、写路径延迟释放的同步机制 \\\tablerowrule
RSS & Resident Set Size & 进程实际驻留物理页，不等于虚拟地址空间大小 \\\tablerowrule
skb & socket buffer & Linux 网络栈中承载包数据和协议元数据的核心对象 \\\tablerowrule
SMP & Symmetric Multiprocessing & 多核并发、per-CPU 数据、runqueue 和 IPI \\\tablerowrule
TLB & Translation Lookaside Buffer & 地址翻译缓存，页表修改后需失效旧项 \\\tablerowrule
VFS & Virtual File System & 统一 superblock、inode、dentry、file 的文件系统抽象 \\\tablerowrule
VMA & Virtual Memory Area & 进程虚拟地址区间策略，用于缺页判断和权限控制 \\\tablerowrule
WAL & Write-Ahead Log & 先写日志再提交数据的崩溃恢复协议 \\
\bottomrule
\end{longtable}
